%==============================================================
%  INTRODUCCIÓN A SQL Y BASES DE DATOS RELACIONALES
%  Módulo 1 – Fundamentos teóricos
%  Formato: Artículo científico / Material académico
%==============================================================

\documentclass[12pt, a4paper]{article}

% ─── Codificación y lenguaje ──────────────────────────────
\usepackage[utf8]{inputenc}
\usepackage[T1]{fontenc}
\usepackage[spanish, es-tabla]{babel}

% ─── Tipografía ───────────────────────────────────────────
\usepackage{lmodern}
\usepackage{microtype}          % Mejora el espaciado tipográfico

% ─── Matemáticas y símbolos ───────────────────────────────
\usepackage{amsmath, amssymb}

% ─── Geometría de página ──────────────────────────────────
\usepackage[
top=2.5cm, bottom=2.5cm,
left=3cm, right=3cm
]{geometry}

% ─── Gráficos y figuras ───────────────────────────────────
\usepackage{graphicx}
\usepackage{float}
\usepackage{caption}
\usepackage{subcaption}
\usepackage[dvipsnames, table]{xcolor}
\usepackage{tikz}

% ─── Tablas ───────────────────────────────────────────────
\usepackage{booktabs}
\usepackage{array}
\usepackage{tabularx}
\usepackage{longtable}
\usepackage{multirow}

% ─── Listados de código ───────────────────────────────────
\usepackage{listings}
\usepackage{lstautogobble}

% Estilo general para código SQL
\lstdefinestyle{sql}{
	language=SQL,
	basicstyle=\ttfamily\small,
	keywordstyle=\bfseries\color{NavyBlue},
	stringstyle=\color{BrickRed},
	commentstyle=\itshape\color{Gray},
	numberstyle=\tiny\color{Gray},
	numbers=left,
	numbersep=8pt,
	stepnumber=1,
	backgroundcolor=\color{gray!6},
	frame=single,
	framerule=0.4pt,
	rulecolor=\color{gray!50},
	tabsize=4,
	showstringspaces=false,
	breaklines=true,
	breakatwhitespace=false,
	captionpos=b,
	autogobble=true,
	xleftmargin=0.5cm,
	morekeywords={
		DATABASE, DATABASES, TABLES, SHOW, USE, ENGINE, CHARSET,
		AUTO_INCREMENT, UNSIGNED, TINYINT, SMALLINT, MEDIUMINT,
		BIGINT, DECIMAL, FLOAT, DOUBLE, DATETIME, TIMESTAMP,
		VARCHAR, TEXT, BLOB, ENUM, SET, DESCRIBE, DESC,
		FOREIGN KEY, REFERENCES, ON DELETE, ON UPDATE,
		CASCADE, SET NULL, RESTRICT
	}
}
\lstset{style=sql}

% ─── Hipervínculos y referencias ──────────────────────────
\usepackage[
colorlinks=true,
linkcolor=NavyBlue,
citecolor=NavyBlue,
urlcolor=NavyBlue
]{hyperref}
\usepackage{cleveref}

% ─── Cabecera y pie de página ─────────────────────────────
\usepackage{fancyhdr}
\pagestyle{fancy}
\fancyhf{}
\fancyhead[L]{\small\textit{Comisión Nacional de Seguros y Fianzas}}
\fancyhead[R]{\small\textit{Jorge Miyoshi Mikami Alonso}}

\fancyfoot[C]{\small\thepage}
\renewcommand{\headrulewidth}{0.4pt}\\


% ─── Cajas informativas ───────────────────────────────────
\usepackage{tcolorbox}
\tcbuselibrary{breakable, skins}

\newtcolorbox{definicion}[1]{
	colback=NavyBlue!5,
	colframe=NavyBlue!60,
	fonttitle=\bfseries\small,
	title=Definición: #1,
	breakable
}

\newtcolorbox{ejemplo}{
	colback=ForestGreen!5,
	colframe=ForestGreen!50,
	fonttitle=\bfseries\small,
	title=Ejemplo práctico,
	breakable
}

\newtcolorbox{nota}{
	colback=Goldenrod!10,
	colframe=Goldenrod!70,
	fonttitle=\bfseries\small,
	title=Nota importante,
	breakable
}

% ─── Enumeraciones y listas ───────────────────────────────
\usepackage{enumitem}
\setlist[itemize]{leftmargin=*, itemsep=2pt, topsep=4pt}
\setlist[enumerate]{leftmargin=*, itemsep=2pt, topsep=4pt}

% ─── Espaciado ────────────────────────────────────────────
\setlength{\parskip}{4pt}
\setlength{\parindent}{0pt}

% ─── Bibliografía ─────────────────────────────────────────
\usepackage[
backend=biber,
style=ieee,
sorting=nyt
]{biblatex}
\addbibresource{referencias.bib}

%==============================================================
%  DATOS DEL DOCUMENTO
%==============================================================
\title{%
	{\large \textsc{Curso de SQL y Bases de Datos Relacionales}}\\[6pt]
	{\LARGE \textbf{SQL}}\\[4pt]
	{\large Inserción, Consulta y Administración de información en bases de datos relacionales basadas en SQL}
}

\author{%
	\textbf{Material de Referencia }\\[2pt]
	\small \texttt{miyoshikamium@gmail.mx}
}

\date{\small Versión 6.0 — \today}

%==============================================================
%  INICIO DEL DOCUMENTO
%==============================================================
\begin{document}
	
	\maketitle
	\thispagestyle{fancy}
	
	%--------------------------------------------------------------
	\begin{abstract}
		\noindent
		Este documento constituye el material teórico de referencia para el
		Módulo 1 del \textit{Curso de SQL y Bases de Datos Relacionales}.
		Se abordan los fundamentos históricos y conceptuales que dieron origen
		a los sistemas de gestión de bases de datos (SGBD) relacionales, la
		evolución del lenguaje SQL desde sus orígenes en los laboratorios de
		IBM hasta los estándares modernos, y los principios del modelo
		relacional propuesto por Edgar F.\ Codd.
		
		El objetivo es proveer al estudiante una base conceptual sólida que
		permita contextualizar cada construcción sintáctica de SQL dentro de
		un marco teórico coherente, facilitando un aprendizaje significativo
		y duradero. El curso progresa desde la consulta básica de datos,
		pasando por la modificación y escritura, hasta la administración
		de esquemas y la integridad referencial.
		
		\smallskip
		\textbf{Palabras clave:} base de datos, modelo relacional, SQL, SGBD,
		álgebra relacional, normalización, ACID, DML, DDL.
	\end{abstract}
	
	\tableofcontents
	\vspace{8pt}
	\hrule
	\vspace{10pt}
	
	%==============================================================
	% PARTE I: FUNDAMENTOS TEÓRICOS
	%==============================================================
	
	\section{El Problema de los Datos: Una Perspectiva Histórica}
	
	Desde los albores de la civilización, la humanidad ha necesitado
	almacenar y recuperar información de manera sistemática.
	Las tablillas de arcilla sumerias (~3000 a.C.), los registros
	contables del Imperio Romano, o los libros de contabilidad de los
	comerciantes medievales son, en esencia, \textit{sistemas de gestión
		de información}: colecciones de datos organizados con un propósito
	definido y un mecanismo de consulta.
	
	La revolución informática de mediados del siglo XX transformó este
	problema milenario en una disciplina formal.
	Con la proliferación de computadoras en organizaciones bancarias,
	gubernamentales y empresariales durante los años 1950 y 1960,
	la necesidad de almacenar grandes volúmenes de datos de forma
	eficiente y recuperarlos con rapidez se volvió crítica.
	
	\subsection{Era pre-relacional (1950–1970)}
	
	Los primeros sistemas de gestión de datos surgieron como soluciones
	ad-hoc a problemas específicos.
	Predominaban dos modelos principales:
	
	\begin{itemize}
		\item \textbf{Modelo jerárquico:} Los datos se organizaban en una
		estructura de árbol. IBM desarrolló IMS
		(\textit{Information Management System}) en 1966 para la misión
		Apolo de la NASA~\cite{date2003introduction}.
		Aunque eficiente para consultas predefinidas, era rígido:
		agregar una nueva relación entre datos requería reestructurar
		físicamente el sistema.
		
		\item \textbf{Modelo de red (CODASYL):} Propuesto por el comité
		CODASYL en 1969, permitía relaciones muchos-a-muchos mediante
		punteros físicos entre registros~\cite{codasyl1971}.
		Era más flexible que el jerárquico, pero extremadamente complejo
		de administrar y altamente dependiente de la estructura física
		del almacenamiento.
	\end{itemize}
	
	El problema fundamental de ambos modelos era la
	\textbf{dependencia físico-lógica}: los programas de aplicación
	debían conocer cómo estaban organizados físicamente los datos en
	disco para poder acceder a ellos. Un cambio en la estructura de
	almacenamiento implicaba reescribir todos los programas.
	
	%==============================================================
	\section{Edgar F.\ Codd y el Modelo Relacional}
	
	En junio de 1970, Edgar Frank Codd —matemático e investigador
	de IBM— publicó el artículo que cambiaría para siempre el mundo de
	las bases de datos:
	\textit{``A Relational Model of Data for Large Shared Data Banks''}
	en la revista \textit{Communications of the ACM}~\cite{codd1970}.
	
	\begin{definicion}{Modelo Relacional Codd, 1970}
		Un modelo de datos en el que la información se representa
		exclusivamente mediante \textbf{relaciones} (tablas bidimensionales),
		y en el que todas las operaciones sobre los datos se expresan
		mediante un conjunto cerrado de operaciones del
		\textbf{álgebra relacional}, sin necesidad de conocer la
		organización física del almacenamiento.
	\end{definicion}
	
	La contribución de Codd fue revolucionaria por tres razones:
	
	\begin{enumerate}
		\item \textbf{Independencia de datos:} Los programas de aplicación
		operan sobre una representación lógica (tablas), completamente
		desacoplada de cómo los datos se almacenan físicamente.
		
		\item \textbf{Fundamento matemático:} El modelo se basa en la
		teoría de conjuntos y la lógica de predicados de primer orden,
		lo que permite razonar formalmente sobre las operaciones.
		
		\item \textbf{Lenguaje declarativo:} El programador especifica
		\emph{qué} datos quiere, no \emph{cómo} obtenerlos.
		El sistema optimiza el acceso automáticamente.
	\end{enumerate}
	
	\subsection{Las 12 reglas de Codd}
	
	En 1985, Codd formalizó su modelo en 12 reglas
	(numeradas del 0 al 12) que debe cumplir un sistema para
	ser considerado verdaderamente relacional~\cite{codd1985}.
	Las más relevantes para este curso son:
	
	\begin{table}[H]
		\centering
		\small
		\renewcommand{\arraystretch}{1.3}
		\begin{tabularx}{\columnwidth}{@{} c X @{}}
			\toprule
			\textbf{Regla} & \textbf{Descripción} \\
			\midrule
			0 & Gestión mediante capacidades relacionales \\
			1 & Representación de la información como tablas \\
			2 & Acceso garantizado a cada dato (tabla + clave + columna) \\
			3 & Tratamiento sistemático de valores nulos \\
			6 & Actualización de vistas \\
			9 & Independencia de datos física y lógica \\
			12 & No subversión del modelo relacional \\
			\bottomrule
		\end{tabularx}
		\caption{Selección de las 12 reglas de Codd (1985).}
		\label{tab:reglas-codd}
	\end{table}
	
	%==============================================================
	\section{Nacimiento de SQL: De SEQUEL a SQL}
	
	A raíz de la publicación de Codd, IBM inició el proyecto
	\textit{System R} en 1973 en sus laboratorios de San José,
	California, con el objetivo de demostrar la viabilidad práctica del
	modelo relacional~\cite{chamberlin1974}.
	Donald D.\ Chamberlin y Raymond F.\ Boyce diseñaron un lenguaje de
	consulta denominado \textbf{SEQUEL}
	(\textit{Structured English Query Language}).
	
	El nombre fue posteriormente acortado a \textbf{SQL}
	(\textit{Structured Query Language}) por razones legales de marca
	registrada.
	En paralelo, Michael Stonebraker y Eugene Wong desarrollaron
	\textbf{INGRES} en la Universidad de California, Berkeley (1973–1977),
	que introdujo su propio lenguaje de consulta llamado QUEL.
	
	\begin{nota}
		Aunque coloquialmente se pronuncia tanto ``sequel'' como
		``ese-cu-ele'', el estándar ISO/IEC lo denomina formalmente
		``SQL''. En este curso usaremos la pronunciación ``ese-cu-ele''
		por su precisión técnica.
	\end{nota}
	
	\subsection{El primer sistema comercial: Oracle}
	
	En 1979, la empresa Relational Software Inc.\ (hoy Oracle
	Corporation), fundada por Larry Ellison, liberó el primer
	sistema de gestión de bases de datos relacionales comercial:
	\textbf{Oracle V2}.
	Paradójicamente, Oracle llegó al mercado antes que el propio
	IBM con su DB2 (1983), a pesar de que la tecnología subyacente
	provenía de los laboratorios de IBM.
	
	%==============================================================
	\section{¿Qué es una Base de Datos?}
	
	\begin{definicion}{Base de Datos}
		Una \textbf{base de datos} es una colección organizada, persistente
		y compartible de datos relacionados entre sí, diseñada para
		minimizar la redundancia y maximizar la integridad, de modo que
		múltiples aplicaciones y usuarios puedan acceder a ella de manera
		concurrente y controlada~\cite{date2003introduction}.
	\end{definicion}
	
	Esta definición merece descomponerse en sus atributos esenciales:
	
	\begin{itemize}
		\item \textbf{Organizada:} Los datos siguen una estructura formal
		(esquema) que define qué tipos de información se almacenan y
		cómo se relacionan entre sí.
		
		\item \textbf{Persistente:} Los datos sobreviven a la terminación
		de los procesos que los crean. A diferencia de las variables en
		memoria, los datos en una base de datos permanecen aunque el
		sistema se apague.
		
		\item \textbf{Compartible:} Múltiples usuarios o aplicaciones
		pueden acceder a los mismos datos simultáneamente, de forma
		controlada y sin interferencias.
		
		\item \textbf{Mínima redundancia:} El mismo dato no se repite
		innecesariamente en múltiples lugares, lo que reduce el riesgo
		de inconsistencias (un dato cambia en un lugar pero no en otro).
	\end{itemize}
	
	\subsection{Analogía con el mundo real}
	
	Imaginemos el sistema de una biblioteca universitaria:
	
	\begin{ejemplo}
		Una biblioteca almacena información sobre \textbf{libros},
		\textbf{estudiantes} y \textbf{préstamos}.
		Sin una base de datos, cada ficha de préstamo copiaría el título
		completo del libro, el autor, el nombre del estudiante, su carrera,
		etc.
		Si un estudiante cambia su dirección, habría que actualizar decenas
		de fichas físicas.
		
		Con una base de datos relacional, el estudiante aparece
		\emph{una sola vez} en la tabla \texttt{Estudiantes}.
		Todos los préstamos simplemente hacen referencia a su identificador.
		Un cambio de dirección se realiza en un único lugar y se refleja
		automáticamente en todo el sistema.
	\end{ejemplo}
	
	%==============================================================
	% PARTE II: ADMINISTRACIÓN DE BASES DE DATOS Y TABLAS
	%==============================================================
	
	\newpage
	\section{Administración de la Base de Datos}
	
	Las sentencias de administración de bases de datos pertenecen al sublenguaje
	\textbf{DDL} (\textit{Data Definition Language}). Operan sobre el contenedor de más alto nivel en MySQL: la
	\textbf{base de datos} (también llamada \textit{schema}), que agrupa tablas,
	vistas, procedimientos y demás objetos.
	
	\subsection{\texttt{CREATE DATABASE}}
	
	\begin{definicion}{CREATE DATABASE}
		Crea una nueva base de datos vacía en el servidor. Permite especificar el
		\textbf{juego de caracteres} (\textit{character set}) y la
		\textbf{colación} (\textit{collation}) que determinarán cómo se almacenan
		y comparan las cadenas de texto en todas sus tablas.
	\end{definicion}
	
	\begin{lstlisting}[style=sql, caption={Sintaxis de CREATE DATABASE}]
		CREATE DATABASE [IF NOT EXISTS] nombre_bd
		[CHARACTER SET  charset_name]
		[COLLATE        collation_name];
	\end{lstlisting}
	
	\begin{ejemplo}
		Crear una base de datos para una aplicación en español con soporte
		completo de Unicode:
		
		\begin{lstlisting}[style=sql, caption={Base de datos con utf8mb4 y collation en español}]
			CREATE DATABASE IF NOT EXISTS tienda_online
			CHARACTER SET  utf8mb4
			COLLATE        utf8mb4_unicode_ci;
			
			-- Seleccionar la base de datos activa
			USE tienda_online;
		\end{lstlisting}
	\end{ejemplo}
	
	\textbf{Juegos de caracteres y colaciones más comunes:}
	
	\begin{table}[H]
		\centering
		\caption{Character sets y collations habituales en MySQL}
		\renewcommand{\arraystretch}{1.3}
		\rowcolors{2}{gray!10}{white}
		\begin{tabularx}{\textwidth}{>{\ttfamily}l >{\ttfamily}l X}
			\toprule
			\textbf{charset}  & \textbf{collation}          & \textbf{Uso recomendado}                        \\
			\midrule
			utf8mb4  & utf8mb4\_unicode\_ci  & Unicode completo; soporta emojis; recomendado   \\
			utf8mb4  & utf8mb4\_0900\_ai\_ci & Unicode 9.0; sin acento/mayúsculas; MySQL 8.0+  \\
			utf8mb4  & utf8mb4\_bin          & Comparación binaria; distingue mayúsculas       \\
			latin1   & latin1\_swedish\_ci   & Legado; solo caracteres ISO-8859-1              \\
			\bottomrule
		\end{tabularx}
	\end{table}
	
	\begin{nota}
		Usar \texttt{utf8} (sin \texttt{mb4}) en MySQL es un error frecuente: ese
		alias solo soporta hasta 3 bytes por carácter y excluye emojis y algunos
		ideogramas CJK. La elección correcta para proyectos nuevos es siempre
		\texttt{utf8mb4}.
	\end{nota}
	
	\begin{lstlisting}[style=sql, caption={Comandos de exploración de bases de datos}]
		SHOW DATABASES;                    -- listar todas las bases de datos
		SHOW CREATE DATABASE tienda_online; -- ver la definición completa
		SELECT DATABASE();                  -- ver la base de datos activa
	\end{lstlisting}
	
	\subsection{\texttt{DROP DATABASE}}
	
	\begin{definicion}{DROP DATABASE}
		Elimina la base de datos y \textbf{todos los objetos que contiene}
		(tablas, vistas, procedimientos, datos) de forma permanente e
		irreversible. No existe confirmación ni papelera de reciclaje.
	\end{definicion}
	
	\begin{lstlisting}[style=sql, caption={Sintaxis de DROP DATABASE}]
		DROP DATABASE [IF EXISTS] nombre_bd;
	\end{lstlisting}
	
	\begin{ejemplo}
		\begin{lstlisting}[style=sql, caption={DROP DATABASE con y sin guarda}]
			-- Genera error si la base de datos no existe
			DROP DATABASE tienda_online;
			
			-- No genera error si no existe (útil en scripts de despliegue)
			DROP DATABASE IF EXISTS tienda_online;
		\end{lstlisting}
	\end{ejemplo}
	
	\begin{nota}
		\textbf{Precauciones antes de ejecutar DROP DATABASE:}
		\begin{enumerate}
			\item Verificar que se tiene un respaldo (\textit{backup}) actualizado.
			\item Confirmar que ninguna aplicación activa tiene conexiones abiertas
			a esa base de datos.
			\item Usar siempre \texttt{IF EXISTS} en scripts de producción.
		\end{enumerate}
	\end{nota}
	
	%==============================================================
	\newpage
	\section{Administración de Tablas}
	
	Las sentencias de administración de tablas pertenecen al sublenguaje
	\textbf{DDL} (\textit{Data Definition Language}). A diferencia de las
	instrucciones DML, los cambios DDL son \textbf{autoconfirmados}
	(\textit{auto-commit}) en la mayoría de los motores: no pueden deshacerse
	con \texttt{ROLLBACK}.
	
	\subsection{\texttt{CREATE TABLE}}
	
	\begin{definicion}{CREATE TABLE}
		Crea una nueva tabla en la base de datos activa. La sentencia define el
		nombre de cada columna, su \textbf{tipo de dato}, las
		\textbf{restricciones} de integridad y las \textbf{opciones de tabla}
		(motor de almacenamiento, juego de caracteres, etc.).
	\end{definicion}
	
	\textbf{Sintaxis completa:}
	
	\begin{lstlisting}[style=sql, caption={Estructura general de CREATE TABLE}]
		CREATE TABLE [IF NOT EXISTS] nombre_tabla (
		columna1  tipo  [opciones_columna] [restriccion_columna],
		columna2  tipo  [opciones_columna] [restriccion_columna],
		...
		[restricciones_tabla]
		) [opciones_tabla];
	\end{lstlisting}
	
	\textbf{Tipos de datos más utilizados en MySQL:}
	
	\begin{table}[H]
		\centering
		\caption{Tipos de datos frecuentes en MySQL}
		\renewcommand{\arraystretch}{1.3}
		\rowcolors{2}{gray!10}{white}
		\begin{tabularx}{\textwidth}{>{\ttfamily}l l X}
			\toprule
			\textbf{Tipo}         & \textbf{Categoría} & \textbf{Descripción}                             \\
			\midrule
			TINYINT               & Entero    & 1 byte; rango $-128$ a $127$ (o $0$–$255$ UNSIGNED) \\
			INT / INTEGER         & Entero    & 4 bytes; rango $\pm2{,}1 \times 10^9$            \\
			BIGINT                & Entero    & 8 bytes; para identificadores de gran volumen    \\
			DECIMAL(p,s)          & Decimal   & Precisión exacta; ideal para valores monetarios  \\
			FLOAT / DOUBLE        & Flotante  & Precisión aproximada; uso científico             \\
			CHAR(n)               & Texto     & Longitud fija $n$; rellena con espacios          \\
			VARCHAR(n)            & Texto     & Longitud variable hasta $n$ caracteres           \\
			TEXT                  & Texto     & Hasta 65\,535 bytes; sin valor DEFAULT           \\
			DATE                  & Fecha     & Formato \texttt{YYYY-MM-DD}                      \\
			DATETIME              & Fecha/hora& Formato \texttt{YYYY-MM-DD HH:MM:SS}             \\
			TIMESTAMP             & Fecha/hora& UTC; se actualiza automáticamente con \texttt{ON UPDATE} \\
			BOOLEAN / TINYINT(1)  & Lógico    & MySQL almacena booleanos como \texttt{TINYINT(1)} \\
			ENUM('a','b',\ldots)  & Catálogo  & Solo acepta los valores de la lista              \\
			\bottomrule
		\end{tabularx}
	\end{table}
	
	\begin{ejemplo}
		Tabla completa con varios tipos de datos, restricciones y opciones:
		
		\begin{lstlisting}[style=sql, caption={CREATE TABLE con opciones reales de producción}]
			CREATE TABLE IF NOT EXISTS productos (
			producto_id   INT UNSIGNED     NOT NULL AUTO_INCREMENT,
			sku           VARCHAR(30)      NOT NULL,
			nombre        VARCHAR(120)     NOT NULL,
			descripcion   TEXT,
			precio        DECIMAL(10,2)    NOT NULL DEFAULT 0.00,
			stock         INT UNSIGNED     NOT NULL DEFAULT 0,
			activo        TINYINT(1)       NOT NULL DEFAULT 1,
			categoria_id  INT UNSIGNED     NOT NULL,
			creado_en     DATETIME         NOT NULL DEFAULT CURRENT_TIMESTAMP,
			actualizado   TIMESTAMP        NOT NULL
			DEFAULT CURRENT_TIMESTAMP
			ON UPDATE CURRENT_TIMESTAMP,
			
			PRIMARY KEY (producto_id),
			UNIQUE  KEY uq_sku (sku),
			CONSTRAINT chk_precio CHECK (precio >= 0),
			CONSTRAINT fk_prod_cat
			FOREIGN KEY (categoria_id)
			REFERENCES categorias(categoria_id)
			ON DELETE RESTRICT
			ON UPDATE CASCADE
			)
			ENGINE  = InnoDB
			CHARSET = utf8mb4
			COLLATE = utf8mb4_unicode_ci;
		\end{lstlisting}
	\end{ejemplo}
	
	\subsection{Restricciones de Integridad}
	
	\subsubsection{\texttt{NOT NULL}}
	
	\begin{definicion}{NOT NULL}
		Prohíbe que una columna almacene el valor \texttt{NULL}. Toda inserción o
		actualización que no suministre un valor para esa columna — y que no tenga
		\texttt{DEFAULT} definido — será rechazada por el motor.
	\end{definicion}
	
	\begin{lstlisting}[style=sql, caption={Columnas con y sin NOT NULL}]
		CREATE TABLE empleados (
		empleado_id  INT          NOT NULL AUTO_INCREMENT,
		nombre       VARCHAR(100) NOT NULL,   -- obligatorio
		telefono     VARCHAR(20)  NULL,       -- opcional (admite NULL)
		correo       VARCHAR(120) NOT NULL,
		PRIMARY KEY (empleado_id)
		);
	\end{lstlisting}
	
	\begin{nota}
		Omitir \texttt{NOT NULL} equivale implícitamente a \texttt{NULL} en MySQL.
		Declararlo de forma explícita mejora la legibilidad y la intención del
		esquema.
	\end{nota}
	
	\subsubsection{\texttt{UNIQUE}}
	
	\begin{definicion}{UNIQUE}
		Garantiza que todos los valores de una columna (o combinación de columnas)
		sean distintos dentro de la tabla. A diferencia de \texttt{PRIMARY KEY},
		una tabla puede tener \textbf{múltiples} restricciones \texttt{UNIQUE} y
		estas \textbf{sí permiten} el valor \texttt{NULL} (con matices según el
		motor).
	\end{definicion}
	
	\begin{lstlisting}[style=sql, caption={UNIQUE a nivel de columna y de tabla}]
		CREATE TABLE usuarios (
		usuario_id  INT          UNSIGNED AUTO_INCREMENT PRIMARY KEY,
		correo      VARCHAR(120) NOT NULL UNIQUE,          -- nivel columna
		username    VARCHAR(50)  NOT NULL,
		doc_tipo    CHAR(3)      NOT NULL,
		doc_numero  VARCHAR(20)  NOT NULL,
		
		-- UNIQUE compuesto a nivel de tabla
		UNIQUE KEY uq_documento (doc_tipo, doc_numero)
		);
	\end{lstlisting}
	
	\begin{table}[H]
		\centering
		\caption{Diferencias entre \texttt{UNIQUE} y \texttt{PRIMARY KEY}}
		\renewcommand{\arraystretch}{1.3}
		\rowcolors{2}{gray!10}{white}
		\begin{tabularx}{\textwidth}{l X X}
			\toprule
			\textbf{Característica} & \textbf{PRIMARY KEY}          & \textbf{UNIQUE}                      \\
			\midrule
			Cantidad por tabla      & Solo una                      & Múltiples                            \\
			Admite NULL             & No                            & Sí (un NULL por columna en MySQL)    \\
			Índice generado         & Clúster (InnoDB)              & Secundario no clúster                \\
			Uso habitual            & Identificador de fila         & Atributos alternativos únicos        \\
			\bottomrule
		\end{tabularx}
	\end{table}
	
	\subsubsection{\texttt{PRIMARY KEY}}
	
	\begin{definicion}{PRIMARY KEY}
		Identifica de forma \textbf{única e irrepetible} cada fila de la tabla.
		Implica \texttt{NOT NULL} y \texttt{UNIQUE} de forma automática. Puede
		ser \textbf{simple} (una columna) o \textbf{compuesta} (varias columnas).
	\end{definicion}
	
	\begin{lstlisting}[style=sql, caption={PK simple vs.\ PK compuesta}]
		-- Clave primaria simple (columna única)
		CREATE TABLE categorias (
		categoria_id  INT UNSIGNED AUTO_INCREMENT,
		nombre        VARCHAR(80) NOT NULL,
		PRIMARY KEY (categoria_id)
		);
		
		-- Clave primaria compuesta (tabla intermedia N:M)
		CREATE TABLE inscripciones (
		estudiante_id  INT UNSIGNED NOT NULL,
		curso_id       INT UNSIGNED NOT NULL,
		fecha_alta     DATE         NOT NULL,
		PRIMARY KEY (estudiante_id, curso_id)   -- par debe ser único
		);
	\end{lstlisting}
	
	\begin{nota}
		En InnoDB, la \texttt{PRIMARY KEY} determina el \textbf{índice clúster}:
		las filas se almacenan físicamente ordenadas por ella. Elegir una PK de
		tipo entero secuencial (\texttt{AUTO\_INCREMENT}) produce inserciones más
		eficientes que una PK de tipo cadena o UUID aleatorio.
	\end{nota}
	
	\subsubsection{\texttt{CHECK}}
	
	\begin{definicion}{CHECK}
		Valida que el valor de una columna cumpla una \textbf{expresión booleana}
		antes de aceptar la inserción o actualización. Si la expresión devuelve
		\texttt{FALSE}, la operación es rechazada. Disponible en MySQL a partir
		de la versión~8.0.16.
	\end{definicion}
	
	\begin{lstlisting}[style=sql, caption={CHECK a nivel de columna y de tabla}]
		CREATE TABLE contratos (
		contrato_id   INT UNSIGNED AUTO_INCREMENT PRIMARY KEY,
		fecha_inicio  DATE         NOT NULL,
		fecha_fin     DATE,
		salario       DECIMAL(10,2) NOT NULL,
		CONSTRAINT chk_fecha CHECK (fecha_fin > fecha_inicio),
		CONSTRAINT chk_salario CHECK (salario > 0)
		);
	\end{lstlisting}
	
	\subsubsection{\texttt{FOREIGN KEY}}
	
	\begin{definicion}{FOREIGN KEY}
		Establece una \textbf{integridad referencial}: la columna o combinación
		de columnas debe coincidir con valores que existen en la clave primaria
		de otra tabla (la tabla referenciada). Las acciones \texttt{ON DELETE}
		y \texttt{ON UPDATE} especifican qué sucede cuando se modifica la
		tabla padre.
	\end{definicion}
	
	\begin{lstlisting}[style=sql, caption={FOREIGN KEY con políticas de actualización}]
		CREATE TABLE pedidos (
		pedido_id   INT UNSIGNED AUTO_INCREMENT PRIMARY KEY,
		cliente_id  INT UNSIGNED NOT NULL,
		fecha       DATE NOT NULL,
		FOREIGN KEY (cliente_id) REFERENCES clientes(cliente_id)
		ON DELETE RESTRICT      -- no permite eliminar si hay pedidos
		ON UPDATE CASCADE       -- actualiza el ID si cambia en clientes
		);
	\end{lstlisting}
	
	\subsubsection{\texttt{DEFAULT}}
	
	\begin{definicion}{DEFAULT}
		Asigna un valor automático a una columna cuando no se proporciona
		uno durante una inserción. Acelera la entrada de datos y reduce
		valores \texttt{NULL} innecesarios.
	\end{definicion}
	
	\begin{lstlisting}[style=sql, caption={DEFAULT con valores estáticos y dinámicos}]
		-- DEFAULT con valores estáticos
		CREATE TABLE usuarios (
		usuario_id   INT UNSIGNED AUTO_INCREMENT PRIMARY KEY,
		activo       TINYINT(1) DEFAULT 1,
		saldo        DECIMAL(10,2) DEFAULT 0.00,
		fecha_registro DATETIME DEFAULT CURRENT_TIMESTAMP
		);
	\end{lstlisting}
	
	\subsubsection{\texttt{AUTO\_INCREMENT}}
	
	\begin{definicion}{AUTO\_INCREMENT}
		Genera automáticamente valores enteros secuenciales para una columna
		en cada nueva inserción. Solo puede haber un \texttt{AUTO\_INCREMENT}
		por tabla y típicamente se usa con \texttt{PRIMARY KEY}.
	\end{definicion}
	
	\begin{lstlisting}[style=sql, caption={AUTO\_INCREMENT en contexto}]
		CREATE TABLE productos (
		producto_id INT UNSIGNED NOT NULL AUTO_INCREMENT,
		nombre      VARCHAR(120) NOT NULL,
		PRIMARY KEY (producto_id)
		);
		
		-- El valor inicial es 1 por defecto
		ALTER TABLE productos AUTO_INCREMENT = 1000;  -- cambiar inicio
	\end{lstlisting}
	
	\begin{nota}
		El contador \texttt{AUTO\_INCREMENT} \textbf{nunca retrocede} al eliminar
		filas. Si se inserta y luego se borra el registro con \texttt{ID = 1005},
		el siguiente ID será~\texttt{1006}, no~\texttt{1005}. Esto es intencional
		para evitar colisiones en entornos concurrentes.
	\end{nota}
	
	\subsection{\texttt{DROP TABLE}}
	
	\begin{definicion}{DROP TABLE}
		Elimina una tabla de forma \textbf{permanente e irreversible}: borra la
		estructura, todos sus datos y los índices asociados. No puede deshacerse
		con \texttt{ROLLBACK}.
	\end{definicion}
	
	\begin{lstlisting}[style=sql, caption={DROP TABLE básico y con guarda}]
		-- Falla con error si la tabla no existe
		DROP TABLE pedidos;
		
		-- No genera error si la tabla no existe
		DROP TABLE IF EXISTS pedidos;
		
		-- Eliminar varias tablas en una sola sentencia
		DROP TABLE IF EXISTS inscripciones, pedidos, sesiones;
	\end{lstlisting}
	
	\textbf{Orden correcto con claves foráneas:}
	
	\begin{lstlisting}[style=sql, caption={DROP TABLE respetando integridad referencial}]
		-- ERROR: no se puede eliminar 'clientes' mientras 'pedidos' la referencia
		DROP TABLE clientes;   -- falla
		
		-- Correcto: eliminar primero la tabla HIJO, luego la tabla PADRE
		DROP TABLE IF EXISTS pedidos;    -- hijo (tiene la FK)
		DROP TABLE IF EXISTS clientes;   -- padre (tiene la PK referenciada)
	\end{lstlisting}
	
	\begin{nota}
		Como alternativa al orden manual, es posible deshabilitar
		temporalmente la comprobación de claves foráneas. Esta práctica debe
		usarse con extrema cautela y solo en scripts de mantenimiento
		controlados:
		\begin{lstlisting}[style=sql, caption={Deshabilitar FK solo para scripts de mantenimiento}]
			SET FOREIGN_KEY_CHECKS = 0;
			DROP TABLE IF EXISTS clientes, pedidos;
			SET FOREIGN_KEY_CHECKS = 1;
		\end{lstlisting}
	\end{nota}
	
	\subsection{\texttt{ALTER TABLE}}
	
	\begin{definicion}{ALTER TABLE}
		Modifica la \textbf{estructura} de una tabla existente sin necesidad de
		recrearla. Permite agregar, renombrar, modificar y eliminar columnas, así
		como gestionar restricciones e índices.
	\end{definicion}
	
	\subsubsection{\texttt{ADD}}
	
	Agrega nuevas columnas o restricciones a una tabla existente.
	
	\begin{lstlisting}[style=sql, caption={ADD COLUMN — agregar columnas}]
		-- Agregar una columna al final (comportamiento por defecto)
		ALTER TABLE productos
		ADD COLUMN descuento DECIMAL(5,2) NOT NULL DEFAULT 0.00;
		
		-- Agregar una columna en una posición específica
		ALTER TABLE productos
		ADD COLUMN imagen_url VARCHAR(255) NULL
		AFTER nombre;
		
		-- Agregar al principio
		ALTER TABLE productos
		ADD COLUMN prioridad TINYINT NOT NULL DEFAULT 0
		FIRST;
	\end{lstlisting}
	
	\begin{lstlisting}[style=sql, caption={ADD CONSTRAINT — agregar restricciones}]
		-- Agregar índice UNIQUE
		ALTER TABLE usuarios
		ADD CONSTRAINT uq_correo UNIQUE (correo);
		
		-- Agregar clave foránea
		ALTER TABLE pedidos
		ADD CONSTRAINT fk_pedido_cliente
		FOREIGN KEY (cliente_id) REFERENCES clientes(cliente_id)
		ON DELETE RESTRICT ON UPDATE CASCADE;
		
		-- Agregar CHECK
		ALTER TABLE productos
		ADD CONSTRAINT chk_descuento CHECK (descuento BETWEEN 0 AND 100);
	\end{lstlisting}
	
	\subsubsection{\texttt{RENAME COLUMN}}
	
	Cambia el nombre de una columna preservando su tipo, restricciones y datos.
	Disponible en MySQL~8.0+.
	
	\begin{lstlisting}[style=sql, caption={RENAME COLUMN}]
		-- Renombrar una columna (MySQL 8.0+)
		ALTER TABLE clientes
		RENAME COLUMN telefono TO telefono_principal;
		
		-- Renombrar varias columnas en una sola sentencia
		ALTER TABLE productos
		RENAME COLUMN img_url   TO imagen_url,
		RENAME COLUMN desc_prod TO descripcion;
	\end{lstlisting}
	
	\begin{nota}
		En versiones anteriores a MySQL~8.0, el renombrado se realiza con
		\texttt{CHANGE COLUMN}, que requiere repetir el tipo de dato:
		\begin{lstlisting}[style=sql, caption={CHANGE COLUMN para MySQL anterior a 8.0}]
			ALTER TABLE clientes
			CHANGE COLUMN telefono telefono_principal VARCHAR(20) NULL;
		\end{lstlisting}
	\end{nota}
	
	\subsubsection{\texttt{MODIFY COLUMN}}
	
	Cambia el \textbf{tipo de dato}, las restricciones o las opciones de una
	columna existente sin alterar su nombre.
	
	\begin{lstlisting}[style=sql, caption={MODIFY COLUMN — cambios de tipo y propiedades}]
		-- Ampliar el tamaño de un VARCHAR
		ALTER TABLE productos
		MODIFY COLUMN nombre VARCHAR(200) NOT NULL;
		
		-- Cambiar tipo y agregar restricción
		ALTER TABLE contratos
		MODIFY COLUMN salario DECIMAL(12,2) NOT NULL DEFAULT 0.00;
		
		-- Cambiar posición de la columna
		ALTER TABLE empleados
		MODIFY COLUMN correo VARCHAR(120) NOT NULL AFTER nombre;
	\end{lstlisting}
	
	\begin{nota}
		\texttt{MODIFY COLUMN} reescribe la definición completa de la columna.
		Omitir \texttt{NOT NULL} o \texttt{DEFAULT} los elimina, aunque el dato
		ya exista. Siempre verificar la definición actual con \texttt{DESCRIBE}
		antes de modificar.
	\end{nota}
	
	\subsubsection{\texttt{DROP COLUMN}}
	
	Elimina una columna y \textbf{todos sus datos} de forma permanente.
	
	\begin{lstlisting}[style=sql, caption={DROP COLUMN}]
		-- Eliminar una columna
		ALTER TABLE productos
		DROP COLUMN imagen_url;
		
		-- Eliminar varias columnas en una sola sentencia
		ALTER TABLE empleados
		DROP COLUMN fax,
		DROP COLUMN extension;
		
		-- Eliminar una restricción antes de poder borrar la columna
		ALTER TABLE pedidos
		DROP FOREIGN KEY fk_pedido_cliente,
		DROP COLUMN cliente_id;
	\end{lstlisting}
	
	\begin{nota}
		Si la columna forma parte de un índice o restricción, MySQL exige
		eliminar primero esa dependencia. Usar \texttt{SHOW CREATE TABLE}
		para identificar los nombres exactos de índices y claves foráneas
		antes de ejecutar el \texttt{DROP}.
	\end{nota}
	
	%==============================================================
	% PARTE III: CONSULTA DE DATOS
	%==============================================================
	
	\newpage
	\section{Consulta de Datos: \texttt{SELECT}}
	
	La instrucción \texttt{SELECT} es, sin duda, la más utilizada y la
	más rica en posibilidades de todo el lenguaje SQL.
	Su propósito es \textbf{consultar} datos: extraer filas de una o
	varias tablas, aplicarles filtros, transformaciones y ordenamientos,
	y devolver un resultado tabular.
	A diferencia de \texttt{INSERT}, \texttt{UPDATE} o \texttt{DELETE},
	\texttt{SELECT} es una operación de \textbf{sólo lectura}: nunca
	modifica los datos almacenados.
	
	\subsection{Sintaxis completa y cláusulas}
	
	La sintaxis formal de \texttt{SELECT} en MySQL/SQL estándar es la
	siguiente.
	Las partes entre corchetes \texttt{[ ]} son opcionales; las palabras
	en mayúsculas son palabras reservadas de SQL:
	
	\begin{lstlisting}[caption={Sintaxis completa de SELECT.}]
		SELECT  [ALL | DISTINCT | DISTINCTROW]
		lista_de_columnas | *
		FROM    nombre_tabla  [[AS] alias_tabla]
		[WHERE      condicion]
		[GROUP BY   columna [ASC | DESC], ...]
		[HAVING     condicion_de_grupo]
		[ORDER BY   columna [ASC | DESC], ...]
		[LIMIT      numero_de_filas  [OFFSET desplazamiento]];
	\end{lstlisting}
	
	Cada cláusula tiene un rol preciso dentro de la consulta.
	La \cref{tab:clausulas-select} resume su función y el orden en que
	el SGBD las \textit{evalúa} internamente (que no siempre coincide
	con el orden en que se escriben):
	
	\begin{table}[H]
		\centering
		\small
		\renewcommand{\arraystretch}{1.4}
		\begin{tabularx}{\linewidth}{@{} l c X @{}}
			\toprule
			\textbf{Cláusula} & \textbf{Orden eval.} & \textbf{Función} \\
			\midrule
			\texttt{FROM}     & 1 & Determina la(s) tabla(s) fuente de datos. \\
			\texttt{WHERE}    & 2 & Filtra filas \textit{antes} de agrupar. \\
			\texttt{GROUP BY} & 3 & Agrupa filas con valores iguales en las columnas indicadas. \\
			\texttt{HAVING}   & 4 & Filtra grupos \textit{después} de agrupar. \\
			\texttt{SELECT}   & 5 & Proyecta (selecciona) las columnas del resultado. \\
			\texttt{DISTINCT} & 6 & Elimina filas duplicadas del resultado proyectado. \\
			\texttt{ORDER BY} & 7 & Ordena el resultado final. \\
			\texttt{LIMIT}    & 8 & Recorta el número de filas devueltas. \\
			\bottomrule
		\end{tabularx}
		\caption{Cláusulas de \texttt{SELECT} y su orden de evaluación interno.}
		\label{tab:clausulas-select}
	\end{table}
	
	\begin{nota}
		El orden de \textbf{escritura} de las cláusulas es fijo y debe
		respetarse exactamente como aparece en la sintaxis de arriba.
		Sin embargo, el SGBD las \textbf{evalúa} en un orden diferente
		(véase \cref{tab:clausulas-select}).
		Esto explica, por ejemplo, por qué no se puede usar un alias
		definido en \texttt{SELECT} dentro de la cláusula \texttt{WHERE}:
		\texttt{WHERE} se evalúa \textit{antes} de que exista el alias.
	\end{nota}
	
	\subsection{Tabla de ejemplo}
	
	Para ilustrar todos los ejemplos de esta sección usaremos una tabla
	llamada \texttt{empleados} con la siguiente estructura y datos:
	
	\begin{lstlisting}[caption={Creación e inicialización de la tabla \texttt{empleados}.}]
		CREATE TABLE empleados (
		id          INT          PRIMARY KEY AUTO_INCREMENT,
		nombre      VARCHAR(60)  NOT NULL,
		departamento VARCHAR(40) NOT NULL,
		cargo       VARCHAR(40)  NOT NULL,
		salario     DECIMAL(10,2) NOT NULL,
		fecha_ingreso DATE        NOT NULL
		);
		
		INSERT INTO empleados (nombre, departamento, cargo, salario, fecha_ingreso) VALUES
		('Ana García',      'Ventas',     'Gerente',    55000.00, '2018-03-15'),
		('Luis Herrera',    'Ventas',     'Analista',   32000.00, '2020-07-01'),
		('Sofía Ramírez',   'TI',         'Desarrolladora', 48000.00, '2019-11-20'),
		('Carlos Mendoza',  'TI',         'DBA',        51000.00, '2017-05-10'),
		('Elena Torres',    'RRHH',       'Coordinadora', 38000.00, '2021-02-28'),
		('Marco Villanueva','Ventas',     'Analista',   31000.00, '2022-09-05'),
		('Patricia Luna',   'TI',         'Desarrolladora', 47000.00, '2020-04-14'),
		('Roberto Salas',   'RRHH',       'Gerente',    52000.00, '2016-08-22'),
		('Diana Castillo',  'Finanzas',   'Analista',   39000.00, '2023-01-10'),
		('Javier Moreno',   'Finanzas',   'Gerente',    58000.00, '2015-06-30');
	\end{lstlisting}
	
	\subsection{Proyección de columnas: \texttt{SELECT *} y columnas específicas}
	
	La \textbf{proyección} es la operación de elegir qué columnas
	aparecerán en el resultado.
	El asterisco \texttt{*} es un comodín que significa
	\textit{``todas las columnas''}.
	
	\begin{lstlisting}[caption={Seleccionar todas las columnas.}]
		-- Devuelve todas las columnas de todos los empleados
		SELECT *
		FROM empleados;
	\end{lstlisting}
	
	\begin{lstlisting}[caption={Seleccionar columnas específicas.}]
		-- Sólo nombre, departamento y salario
		SELECT nombre, departamento, salario
		FROM empleados;
	\end{lstlisting}
	
	\begin{nota}
		En producción se recomienda \textbf{evitar} \texttt{SELECT *} porque:
		(1) transfiere columnas innecesarias por la red,
		(2) si alguien agrega una columna sensible a la tabla, la consulta
		la expondrá automáticamente,
		(3) dificulta la legibilidad del código.
		Úsalo sólo para exploración interactiva.
	\end{nota}
	
	También es posible incluir \textbf{expresiones} en la lista de
	columnas: literales, operaciones aritméticas o llamadas a funciones:
	
	\begin{lstlisting}[caption={Columnas calculadas y literales en SELECT.}]
		SELECT
		nombre,
		salario,
		salario * 12          AS salario_anual,
		salario * 0.10        AS bono_estimado,
		'Activo'              AS estado   -- literal de texto
		FROM empleados;
	\end{lstlisting}
	
	\subsection{Operadores aritméticos}
	
	Dentro de \texttt{SELECT} (y también en \texttt{WHERE}) se pueden
	usar operadores aritméticos estándar sobre columnas numéricas:
	
	\begin{table}[H]
		\centering
		\small
		\renewcommand{\arraystretch}{1.4}
		\begin{tabularx}{\linewidth}{@{} c l X @{}}
			\toprule
			\textbf{Op.} & \textbf{Nombre} & \textbf{Ejemplo} \\
			\midrule
			\texttt{+}  & Suma            & \texttt{salario + 5000} \\
			\texttt{-}  & Resta           & \texttt{salario - descuento} \\
			\texttt{*}  & Multiplicación  & \texttt{salario * 1.15} \\
			\texttt{/}  & División        & \texttt{salario / 14} \\
			\texttt{\%} & Módulo (resto)  & \texttt{id \% 2} (par/impar) \\
			\bottomrule
		\end{tabularx}
		\caption{Operadores aritméticos en SQL.}
		\label{tab:operadores-aritmeticos}
	\end{table}
	
	\begin{lstlisting}[caption={Uso de operadores aritméticos.}]
		SELECT
		nombre,
		salario,
		salario * 1.15        AS salario_con_aumento,
		(salario / 30)        AS salario_diario,
		salario - 5000        AS salario_base
		FROM empleados
		WHERE departamento = 'TI';
	\end{lstlisting}
	
	\subsection{Operadores de comparación}
	
	Los operadores de comparación se usan principalmente en la cláusula
	\texttt{WHERE} para filtrar filas.
	Devuelven un valor booleano (\texttt{TRUE} o \texttt{FALSE}):
	
	\begin{table}[H]
		\centering
		\small
		\renewcommand{\arraystretch}{1.4}
		\begin{tabularx}{\linewidth}{@{} c l X @{}}
			\toprule
			\textbf{Operador} & \textbf{Significado} & \textbf{Ejemplo} \\
			\midrule
			\texttt{=}    & Igual a              & \texttt{departamento = 'Ventas'} \\
			\texttt{<>} o \texttt{!=} & Distinto de & \texttt{cargo <> 'Gerente'} \\
			\texttt{>}    & Mayor que            & \texttt{salario > 40000} \\
			\texttt{<}    & Menor que            & \texttt{salario < 35000} \\
			\texttt{>=}   & Mayor o igual que    & \texttt{salario >= 50000} \\
			\texttt{<=}   & Menor o igual que    & \texttt{salario <= 48000} \\
			\bottomrule
		\end{tabularx}
		\caption{Operadores de comparación en SQL.}
		\label{tab:operadores-comparacion}
	\end{table}
	
	\begin{lstlisting}[caption={Filtros con operadores de comparación.}]
		-- Empleados con salario mayor a 45,000
		SELECT nombre, cargo, salario
		FROM   empleados
		WHERE  salario > 45000;
		
		-- Empleados que NO son del departamento de TI
		SELECT nombre, departamento
		FROM   empleados
		WHERE  departamento <> 'TI';
		
		-- Empleados con salario entre 30,000 y 40,000 (inclusive)
		SELECT nombre, salario
		FROM   empleados
		WHERE  salario >= 30000
		AND  salario <= 40000;
	\end{lstlisting}
	
	%==============================================================
	\section{Modificadores: Filtrado, Ordenamiento y Agrupación}
	
	Los \textbf{modificadores} son cláusulas y operadores que se
	combinan con \texttt{SELECT} para refinar, filtrar y ordenar los
	resultados de una consulta.
	En esta sección cubrimos los modificadores fundamentales que
	todo programador SQL debe dominar.
	
	\subsection{\texttt{DISTINCT}}
	
	Por defecto, \texttt{SELECT} devuelve \textbf{todas} las filas que
	satisfacen la consulta, incluyendo duplicados.
	La palabra clave \texttt{DISTINCT} instruye al SGBD a eliminar las
	filas repetidas del resultado, conservando sólo valores únicos.
	
	\begin{definicion}{\texttt{DISTINCT}}
		Modificador que se coloca inmediatamente después de \texttt{SELECT}
		y elimina del resultado todas las filas cuyo conjunto de valores en
		las columnas proyectadas sea idéntico a otra fila ya presente.
	\end{definicion}
	
	\begin{lstlisting}[caption={DISTINCT sobre una columna.}]
		-- ¿Qué departamentos existen en la empresa?
		-- Sin DISTINCT aparecerían 10 filas (una por empleado)
		SELECT DISTINCT departamento
		FROM   empleados;
	\end{lstlisting}
	
	\begin{lstlisting}[caption={DISTINCT sobre múltiples columnas.}]
		-- Combinaciones únicas de departamento + cargo
		SELECT DISTINCT departamento, cargo
		FROM   empleados
		ORDER BY departamento, cargo;
	\end{lstlisting}
	
	\begin{nota}
		\texttt{DISTINCT} opera sobre la \textit{combinación completa} de
		todas las columnas proyectadas, no columna por columna.
	\end{nota}
	
	\subsection{\texttt{WHERE}}
	
	La cláusula \texttt{WHERE} es el mecanismo de \textbf{filtrado por
		fila} de SQL.
	Evalúa una condición booleana para cada fila de la tabla fuente;
	sólo las filas donde la condición resulta \texttt{TRUE} pasan al
	resultado.
	
	\begin{definicion}{\texttt{WHERE}}
		Cláusula que restringe las filas devueltas por una consulta a
		aquellas que satisfacen una expresión booleana dada.
		Se evalúa \textit{antes} de \texttt{GROUP BY}, \texttt{HAVING} y
		\texttt{SELECT}.
	\end{definicion}
	
	\begin{lstlisting}[caption={Filtros básicos con WHERE.}]
		-- Empleados del departamento de TI
		SELECT nombre, cargo, salario
		FROM   empleados
		WHERE  departamento = 'TI';
		
		-- Empleados con salario mayor a 45,000
		SELECT nombre, salario
		FROM   empleados
		WHERE  salario > 45000;
		
		-- Empleados que ingresaron después del 1 de enero de 2020
		SELECT nombre, fecha_ingreso
		FROM   empleados
		WHERE  fecha_ingreso > '2020-01-01';
	\end{lstlisting}
	
	\subsection{\texttt{ORDER BY}}
	
	\texttt{ORDER BY} define el \textbf{criterio de ordenamiento} del
	resultado final.
	Sin él, el SGBD no garantiza ningún orden particular entre
	ejecuciones.
	
	\begin{lstlisting}[caption={Orden ascendente y descendente.}]
		-- Orden ascendente (A → Z, menor → mayor): por defecto
		SELECT nombre, salario
		FROM   empleados
		ORDER BY salario ASC;   -- ASC es opcional, es el valor por defecto
		
		-- Orden descendente (mayor → menor)
		SELECT nombre, salario
		FROM   empleados
		ORDER BY salario DESC;
	\end{lstlisting}
	
	\begin{lstlisting}[caption={Ordenamiento compuesto por múltiples columnas.}]
		-- Primero por departamento (A→Z),
		-- dentro de cada departamento por salario (mayor→menor)
		SELECT nombre, departamento, salario
		FROM   empleados
		ORDER BY departamento ASC,
		salario      DESC;
	\end{lstlisting}
	
	\subsection{\texttt{LIKE}}
	
	\texttt{LIKE} permite comparar una columna de texto contra un
	\textbf{patrón}, usando caracteres comodín.
	Es la herramienta de búsqueda de texto parcial más básica de SQL.
	
	\begin{table}[H]
		\centering
		\small
		\renewcommand{\arraystretch}{1.4}
		\begin{tabularx}{\linewidth}{@{} c l X @{}}
			\toprule
			\textbf{Comodín} & \textbf{Significado} & \textbf{Ejemplo} \\
			\midrule
			\texttt{\%} & Cero o más caracteres cualesquiera
			& \texttt{'Ana\%'} coincide con \texttt{Ana}, \texttt{Analista}, \texttt{Anabel} \\
			\texttt{\_} & Exactamente un carácter cualquiera
			& \texttt{'\_na'} coincide con \texttt{Ana}, \texttt{Ina}, \texttt{Una} \\
			\bottomrule
		\end{tabularx}
		\caption{Comodines del operador \texttt{LIKE}.}
		\label{tab:comodines-like}
	\end{table}
	
	\begin{lstlisting}[caption={Patrones con LIKE.}]
		-- Nombres que empiezan con 'A'
		SELECT nombre FROM empleados
		WHERE  nombre LIKE 'A%';
		
		-- Nombres que terminan con 'ez'
		SELECT nombre FROM empleados
		WHERE  nombre LIKE '%ez';
		
		-- Nombres que contienen 'ar' en cualquier posición
		SELECT nombre FROM empleados
		WHERE  nombre LIKE '%ar%';
		
		-- Nombres que NO contienen 'a' (NOT LIKE)
		SELECT nombre FROM empleados
		WHERE  nombre NOT LIKE '%a%';
	\end{lstlisting}
	
	\begin{nota}
		En MySQL, \texttt{LIKE} es \textbf{insensible a mayúsculas/minúsculas}
		por defecto cuando la columna usa un \textit{collation} insensible.
		Si necesitas distinguir mayúsculas usa \texttt{LIKE BINARY}.
	\end{nota}
	
	\subsection{\texttt{AND}, \texttt{OR}, \texttt{NOT}}
	
	Los operadores lógicos permiten \textbf{combinar múltiples
		condiciones} dentro de una cláusula \texttt{WHERE}.
	
	\begin{table}[H]
		\centering
		\small
		\renewcommand{\arraystretch}{1.5}
		\begin{tabularx}{\linewidth}{@{} l X X @{}}
			\toprule
			\textbf{Operador} & \textbf{Resultado TRUE cuando\ldots} & \textbf{Precedencia} \\
			\midrule
			\texttt{NOT} & La condición que sigue es \texttt{FALSE}
			& Mayor (se evalúa primero) \\
			\texttt{AND} & \textbf{Ambas} condiciones son \texttt{TRUE}
			& Media \\
			\texttt{OR}  & \textbf{Al menos una} condición es \texttt{TRUE}
			& Menor (se evalúa al final) \\
			\bottomrule
		\end{tabularx}
		\caption{Operadores lógicos y su precedencia.}
		\label{tab:operadores-logicos}
	\end{table}
	
	\begin{lstlisting}[caption={Uso de AND, OR y NOT.}]
		-- AND: empleados de TI con salario mayor a 47,000
		SELECT nombre, departamento, salario
		FROM   empleados
		WHERE  departamento = 'TI'
		AND  salario > 47000;
		
		-- OR: empleados de Ventas O de Finanzas
		SELECT nombre, departamento
		FROM   empleados
		WHERE  departamento = 'Ventas'
		OR   departamento = 'Finanzas';
		
		-- NOT: empleados que NO son gerentes
		SELECT nombre, cargo
		FROM   empleados
		WHERE  NOT cargo = 'Gerente';
		-- equivalente: WHERE cargo <> 'Gerente'
	\end{lstlisting}
	
	\subsection{\texttt{IN} y \texttt{BETWEEN}}
	
	\begin{lstlisting}[caption={IN para búsquedas en conjuntos.}]
		-- Empleados en los departamentos Ventas, TI o Finanzas
		SELECT nombre, departamento
		FROM   empleados
		WHERE  departamento IN ('Ventas', 'TI', 'Finanzas');
		
		-- Equivalente usando OR (menos legible)
		WHERE  departamento = 'Ventas'
		OR departamento = 'TI'
		OR departamento = 'Finanzas';
	\end{lstlisting}
	
	\begin{lstlisting}[caption={BETWEEN para rangos.}]
		-- Empleados con salario entre 40,000 y 50,000 (inclusive)
		SELECT nombre, salario
		FROM   empleados
		WHERE  salario BETWEEN 40000 AND 50000;
		
		-- Equivalente usando AND (menos legible)
		WHERE  salario >= 40000
		AND  salario <= 50000;
	\end{lstlisting}
	
	\subsection{\texttt{GROUP BY} y \texttt{HAVING}}
	
	\texttt{GROUP BY} divide las filas en grupos según los valores
	de una o más columnas. Típicamente se combina con funciones de
	\textbf{agregación} (\texttt{COUNT}, \texttt{SUM}, \texttt{AVG},
	\texttt{MIN}, \texttt{MAX}) para calcular resúmenes por grupo.
	
	\begin{definicion}{GROUP BY}
		Cláusula que divide las filas del resultado de \texttt{WHERE} en
		grupos, donde cada grupo comparte el mismo valor en las columnas
		de agrupamiento.
		Cada grupo produce exactamente \textbf{una fila} en el resultado.
	\end{definicion}
	
	\begin{lstlisting}[caption={GROUP BY con funciones de agregación.}]
		-- Número de empleados y salario promedio por departamento
		SELECT departamento,
		COUNT(*)           AS num_empleados,
		ROUND(AVG(salario), 2) AS promedio,
		SUM(salario)       AS masa_salarial,
		MIN(salario)       AS minimo,
		MAX(salario)       AS maximo
		FROM   empleados
		GROUP BY departamento
		ORDER BY masa_salarial DESC;
	\end{lstlisting}
	
	\texttt{HAVING} filtra \textbf{grupos} (no filas individuales).
	Es el equivalente de \texttt{WHERE} pero aplicado después de que
	\texttt{GROUP BY} ha formado los grupos.
	
	\begin{definicion}{\texttt{HAVING}}
		Cláusula que filtra los grupos producidos por \texttt{GROUP BY},
		conservando sólo aquellos cuya condición (generalmente sobre
		una función de agregación) es \texttt{TRUE}.
	\end{definicion}
	
	\begin{lstlisting}[caption={HAVING: filtrar grupos por condición agregada.}]
		-- Departamentos donde el salario promedio supera 45,000
		SELECT departamento,
		ROUND(AVG(salario), 2) AS promedio
		FROM   empleados
		GROUP BY departamento
		HAVING AVG(salario) > 45000;
		
		-- Departamentos con más de 2 empleados
		SELECT departamento,
		COUNT(*) AS num_empleados
		FROM   empleados
		GROUP BY departamento
		HAVING COUNT(*) > 2;
	\end{lstlisting}
	
	\begin{nota}
		La diferencia entre \texttt{WHERE} y \texttt{HAVING}:
		\texttt{WHERE} descarta filas \textit{antes} de agrupar
		(actúa sobre datos individuales);
		\texttt{HAVING} descarta grupos \textit{después} de agrupar
		(actúa sobre resultados de agregación).
	\end{nota}
	
	\subsection{\texttt{CASE}}
	
	\texttt{CASE} es la expresión de \textbf{lógica condicional} de SQL:
	permite evaluar condiciones dentro de una consulta y devolver
	valores diferentes según cuál se cumpla.
	
	\begin{lstlisting}[caption={CASE buscado (evalúa condiciones booleanas).}]
		-- Clasifica el salario en bandas
		SELECT nombre,
		salario,
		CASE
		WHEN salario >= 55000 THEN 'Banda A (Senior)'
		WHEN salario >= 45000 THEN 'Banda B (Mid)'
		WHEN salario >= 35000 THEN 'Banda C (Junior)'
		ELSE                       'Banda D (Entrada)'
		END AS banda_salarial
		FROM   empleados
		ORDER BY salario DESC;
	\end{lstlisting}
	
	\subsection{\texttt{COALESCE} e \texttt{IFNULL}}
	
	Ambas funciones permiten sustituir un valor \texttt{NULL} por un
	valor alternativo.
	
	\begin{lstlisting}[caption={COALESCE: devuelve el primer valor no NULL.}]
		-- Devuelve el primer contacto disponible entre tres opciones
		SELECT nombre,
		COALESCE(telefono_celular,
		telefono_fijo,
		email,
		'Sin contacto') AS contacto_preferido
		FROM   empleados;
	\end{lstlisting}
	
	\begin{nota}
		\texttt{IFNULL(a, b)} es equivalente a \texttt{COALESCE(a, b)},
		pero \texttt{COALESCE} acepta cualquier número de argumentos y
		es parte del estándar SQL.
		Prefiere \texttt{COALESCE} para mayor portabilidad.
	\end{nota}
	
	\subsection{Combinación de tablas: \texttt{JOIN}}
	
	El operador \texttt{JOIN} combina filas de dos o más tablas basándose
	en una condición de asociación (típicamente igualdad de clave foránea
	y clave primaria).
	
	\subsubsection{\texttt{INNER JOIN}}
	
	\begin{definicion}{INNER JOIN}
		Devuelve solo las filas donde existe una coincidencia exacta en ambas
		tablas. Si una fila de la tabla izquierda no tiene par en la tabla
		derecha, se descarta completamente.
	\end{definicion}
	
	\textbf{Sintaxis general:}
	
	\begin{lstlisting}[style=sql, caption={Sintaxis de INNER JOIN}]
		SELECT columnas
		FROM tabla_izquierda
		INNER JOIN tabla_derecha
		ON tabla_izquierda.clave = tabla_derecha.clave;
	\end{lstlisting}
	
	\begin{ejemplo}
		Listar clientes que han realizado pedidos:
		
		\begin{lstlisting}[style=sql, caption={INNER JOIN: clientes con pedidos}]
			SELECT c.nombre,
			p.pedido_id,
			p.monto
			FROM clientes AS c
			INNER JOIN pedidos AS p
			ON c.cliente_id = p.cliente_id;
		\end{lstlisting}
		
		Solo aparecen clientes que han realizado al menos un pedido.
	\end{ejemplo}
	
	\begin{nota}
		La palabra clave \texttt{INNER} es opcional; \texttt{JOIN} a secas equivale a
		\texttt{INNER JOIN} en todos los motores SQL principales.
	\end{nota}
	
	\subsubsection{\texttt{LEFT JOIN}}
	
	\begin{definicion}{LEFT JOIN}
		Devuelve \textbf{todas las filas de la tabla izquierda} más las coincidencias de la tabla
		derecha. Cuando no existe par, las columnas de la tabla derecha se rellenan con
		\texttt{NULL}.
	\end{definicion}
	
	\begin{lstlisting}[style=sql, caption={Sintaxis de LEFT JOIN}]
		SELECT columnas
		FROM tabla_izquierda
		LEFT JOIN tabla_derecha
		ON tabla_izquierda.clave = tabla_derecha.clave;
	\end{lstlisting}
	
	\begin{ejemplo}
		Listar todos los clientes y, si tienen pedidos, mostrar su monto:
		
		\begin{lstlisting}[style=sql, caption={LEFT JOIN: todos los clientes con sus pedidos}]
			SELECT c.nombre,
			p.pedido_id,
			p.monto
			FROM clientes AS c
			LEFT JOIN pedidos AS p
			ON c.cliente_id = p.cliente_id;
		\end{lstlisting}
		
		Los clientes sin pedidos aparecen con \texttt{NULL} en las columnas de \texttt{pedidos}.
	\end{ejemplo}
	
	\textbf{Caso de uso habitual — detectar registros sin par:}
	
	\begin{lstlisting}[style=sql, caption={Clientes que nunca han realizado un pedido}]
		SELECT c.nombre
		FROM clientes AS c
		LEFT JOIN pedidos AS p
		ON c.cliente_id = p.cliente_id
		WHERE p.pedido_id IS NULL;
	\end{lstlisting}
	
	\subsubsection{\texttt{RIGHT JOIN}}
	
	\begin{definicion}{RIGHT JOIN}
		Devuelve \textbf{todas las filas de la tabla derecha} más las coincidencias de la tabla
		izquierda. Es el simétrico de \texttt{LEFT JOIN}.
	\end{definicion}
	
	\begin{nota}
		En la práctica, \texttt{RIGHT JOIN} se usa con menos frecuencia que \texttt{LEFT JOIN}.
		Cualquier consulta con \texttt{RIGHT JOIN} puede reescribirse como un \texttt{LEFT JOIN}
		intercambiando el orden de las tablas:
		\begin{center}
			\texttt{A RIGHT JOIN B} $\equiv$ \texttt{B LEFT JOIN A}
		\end{center}
	\end{nota}
	
	\subsubsection{\texttt{UNION}}
	
	\begin{definicion}{UNION}
		Combina verticalmente los resultados de dos o más sentencias \texttt{SELECT}.
		\texttt{UNION} elimina filas duplicadas; \texttt{UNION ALL} las conserva.
	\end{definicion}
	
	\begin{lstlisting}[style=sql, caption={Sintaxis de UNION / UNION ALL}]
		SELECT columna1, columna2, ...
		FROM tabla_a
		UNION [ALL]
		SELECT columna1, columna2, ...
		FROM tabla_b;
	\end{lstlisting}
	
	\begin{ejemplo}
		Obtener una lista unificada de correos de dos tablas distintas:
		
		\begin{lstlisting}[style=sql, caption={UNION de correos de clientes y proveedores}]
			SELECT nombre, correo, 'cliente'   AS origen
			FROM clientes
			UNION
			SELECT nombre, correo, 'proveedor' AS origen
			FROM proveedores
			ORDER BY nombre;
		\end{lstlisting}
	\end{ejemplo}
	
	%==============================================================
	% PARTE IV: ESCRITURA Y MODIFICACIÓN DE DATOS
	%==============================================================
	
	\newpage
	\section{Escritura de Datos}
	
	Las sentencias de escritura forman parte del sublenguaje \textbf{DML}
	(\textit{Data Manipulation Language}) de SQL. Permiten insertar, modificar y
	eliminar filas en las tablas de una base de datos relacional.
	
	\begin{nota}
		Toda operación DML es transaccional. En motores como MySQL/InnoDB o PostgreSQL,
		los cambios pueden deshacerse con \texttt{ROLLBACK} mientras la transacción
		no haya sido confirmada con \texttt{COMMIT}.
	\end{nota}
	
	\subsection{\texttt{INSERT}}
	
	\begin{definicion}{INSERT}
		Agrega \textbf{una o más filas} nuevas a una tabla. Los valores proporcionados
		deben respetar los tipos de datos, restricciones \texttt{NOT NULL} y claves
		foráneas definidas en el esquema.
	\end{definicion}
	
	\textbf{Sintaxis — una sola fila:}
	
	\begin{lstlisting}[style=sql, caption={INSERT de una fila con columnas explícitas}]
		INSERT INTO nombre_tabla (col1, col2, col3)
		VALUES (val1, val2, val3);
	\end{lstlisting}
	
	\textbf{Sintaxis — múltiples filas:}
	
	\begin{lstlisting}[style=sql, caption={INSERT de varias filas en una sola sentencia}]
		INSERT INTO nombre_tabla (col1, col2, col3)
		VALUES
		(val1a, val2a, val3a),
		(val1b, val2b, val3b),
		(val1c, val2c, val3c);
	\end{lstlisting}
	
	\begin{ejemplo}
		Poblar la tabla \texttt{productos} con tres artículos:
		
		\begin{lstlisting}[style=sql, caption={Inserción múltiple en la tabla productos}]
			INSERT INTO productos (nombre, precio, stock, categoria_id)
			VALUES
			('Teclado mecánico',  850.00, 40, 3),
			('Monitor 24"',      3200.00, 15, 3),
			('Silla ergonómica', 5400.00,  8, 5);
		\end{lstlisting}
	\end{ejemplo}
	
	\textbf{Inserción a partir de un \texttt{SELECT}:}
	
	\begin{lstlisting}[style=sql, caption={INSERT ... SELECT: copiar filas desde otra tabla}]
		INSERT INTO productos_archivo (nombre, precio, categoria_id)
		SELECT nombre, precio, categoria_id
		FROM   productos
		WHERE  stock = 0;
	\end{lstlisting}
	
	\subsection{\texttt{UPDATE}}
	
	\begin{definicion}{UPDATE}
		Modifica los valores de \textbf{una o más columnas} en las filas que satisfacen
		la condición \texttt{WHERE}. Si se omite \texttt{WHERE}, la actualización afecta
		a \textbf{todas} las filas de la tabla.
	\end{definicion}
	
	\textbf{Sintaxis general:}
	
	\begin{lstlisting}[style=sql, caption={Sintaxis de UPDATE con WHERE}]
		UPDATE nombre_tabla
		SET    col1 = nuevo_val1,
		col2 = nuevo_val2
		WHERE  condicion;
	\end{lstlisting}
	
	\begin{ejemplo}
		Aplicar un incremento de 10\,\% al precio de todos los productos de la
		categoría 3:
		
		\begin{lstlisting}[style=sql, caption={UPDATE con expresión aritmética}]
			UPDATE productos
			SET    precio = precio * 1.10
			WHERE  categoria_id = 3;
		\end{lstlisting}
	\end{ejemplo}
	
	\begin{nota}
		\textbf{Buenas prácticas antes de ejecutar un UPDATE:}
		\begin{enumerate}
			\item Verificar la condición con un \texttt{SELECT} previo usando el mismo
			\texttt{WHERE}.
			\item Ejecutar dentro de una transacción explícita para poder hacer
			\texttt{ROLLBACK} si el resultado no es el esperado.
			\item En MySQL, activar \texttt{sql\_safe\_updates = 1} para bloquear
			\texttt{UPDATE} sin \texttt{WHERE}.
		\end{enumerate}
	\end{nota}
	
	\subsection{\texttt{DELETE}}
	
	\begin{definicion}{DELETE}
		Elimina las filas que cumplen la condición \texttt{WHERE}. Al igual que
		\texttt{UPDATE}, si se omite \texttt{WHERE} se eliminan \textbf{todas} las
		filas de la tabla.
	\end{definicion}
	
	\textbf{Sintaxis general:}
	
	\begin{lstlisting}[style=sql, caption={Sintaxis de DELETE con WHERE}]
		DELETE FROM nombre_tabla
		WHERE condicion;
	\end{lstlisting}
	
	\begin{ejemplo}
		Eliminar los productos sin stock de la categoría 5:
		
		\begin{lstlisting}[style=sql, caption={DELETE con condición compuesta}]
			DELETE FROM productos
			WHERE stock = 0
			AND categoria_id = 5;
		\end{lstlisting}
	\end{ejemplo}
	
	\subsubsection{Diferencias entre \texttt{DELETE} y \texttt{TRUNCATE}}
	
	\begin{table}[H]
		\centering
		\caption{Comparación entre \texttt{DELETE} y \texttt{TRUNCATE}}
		\renewcommand{\arraystretch}{1.3}
		\rowcolors{2}{gray!10}{white}
		\begin{tabularx}{\textwidth}{>{\bfseries}l X X}
			\toprule
			\textbf{Característica}     & \textbf{\texttt{DELETE}}                        & \textbf{\texttt{TRUNCATE}}                   \\
			\midrule
			Sublenguaje SQL             & DML                                             & DDL                                          \\
			Cláusula \texttt{WHERE}     & Sí — elimina filas selectivamente               & No — elimina todas las filas                 \\
			Transaccional               & Sí — admite \texttt{ROLLBACK}                   & Depende del motor; en MySQL no               \\
			Velocidad                   & Más lenta (registro fila a fila)                & Más rápida (libera páginas de datos)         \\
			Reinicio de \texttt{AUTO\_INCREMENT} & No — el contador no se reinicia       & Sí — el contador vuelve a 1                  \\
			\bottomrule
		\end{tabularx}
	\end{table}
	
	\begin{nota}
		Ante la duda entre \texttt{DELETE} y \texttt{TRUNCATE}, preferir
		\texttt{DELETE} dentro de una transacción explícita: ofrece mayor control y
		trazabilidad.
	\end{nota}
	
	%==============================================================
	\newpage
	\section{Almacenamiento de Datos Relacionados}
	
	Cuando una base de datos posee \textbf{integridad referencial} activa, el orden
	y la forma en que se insertan los registros importan tanto como el contenido
	mismo. Violar el orden de inserción provoca errores de clave foránea que el
	motor rechaza antes de escribir cualquier dato.
	
	\begin{nota}
		Regla general: insertar siempre primero la tabla \textbf{padre} (la que
		contiene la clave primaria referenciada) y después la tabla \textbf{hijo}
		(la que contiene la clave foránea). Para eliminar, el orden es el inverso.
	\end{nota}
	
	\subsection{Inserción en Tablas con Relación 1:1}
	
	\begin{definicion}{Relación 1:1}
		Una fila de la tabla \textbf{A} se corresponde con \textbf{exactamente una}
		fila de la tabla \textbf{B}, y viceversa. Se implementa colocando la clave
		foránea en la tabla secundaria con restricción \texttt{UNIQUE}.
	\end{definicion}
	
	\textbf{Esquema de ejemplo:}
	
	\begin{lstlisting}[style=sql, caption={Tablas con relación 1:1 — empleado y expediente}]
		CREATE TABLE empleados (
		empleado_id   INT          UNSIGNED AUTO_INCREMENT PRIMARY KEY,
		nombre        VARCHAR(100) NOT NULL,
		departamento  VARCHAR(60)  NOT NULL
		);
		
		CREATE TABLE expedientes (
		expediente_id INT          UNSIGNED AUTO_INCREMENT PRIMARY KEY,
		empleado_id   INT          UNSIGNED NOT NULL UNIQUE,   -- UNIQUE fuerza 1:1
		nss           CHAR(11)     NOT NULL,
		fecha_alta    DATE         NOT NULL,
		FOREIGN KEY (empleado_id) REFERENCES empleados(empleado_id)
		ON DELETE CASCADE
		ON UPDATE CASCADE
		);
	\end{lstlisting}
	
	\begin{lstlisting}[style=sql, caption={Inserción respetando integridad referencial 1:1}]
		-- Paso 1: insertar en la tabla PADRE
		INSERT INTO empleados (nombre, departamento)
		VALUES ('Laura Mendoza', 'Ingeniería'),
		('Roberto Solis', 'Finanzas');
		
		-- Paso 2: insertar en la tabla HIJO referenciando los IDs existentes
		INSERT INTO expedientes (empleado_id, nss, fecha_alta)
		VALUES (LAST_INSERT_ID(), '578-12-3456', '2021-03-15');
	\end{lstlisting}
	
	\subsection{Inserción en Tablas con Relación 1:N}
	
	\begin{definicion}{Relación 1:N}
		Una fila de la tabla \textbf{padre} puede asociarse con \textbf{múltiples}
		filas de la tabla \textbf{hijo}, pero cada fila hijo pertenece a un único
		padre. Es la cardinalidad más frecuente en bases de datos relacionales.
	\end{definicion}
	
	\begin{lstlisting}[style=sql, caption={Inserción 1:N — un cliente con varios pedidos}]
		-- Paso 1: insertar el cliente (tabla padre)
		INSERT INTO clientes (nombre, correo)
		VALUES ('Ana Torres', 'ana.torres@correo.com');
		
		-- Paso 2: insertar sus pedidos (tabla hijo)
		INSERT INTO pedidos (cliente_id, fecha, total)
		VALUES
		(LAST_INSERT_ID(), '2024-01-15', 1250.00),
		(LAST_INSERT_ID(), '2024-03-02',  340.50),
		(LAST_INSERT_ID(), '2024-06-18', 2890.75);
	\end{lstlisting}
	
	\subsection{Inserción en Tablas con Relación N:M}
	
	\begin{definicion}{Relación N:M}
		Múltiples filas de la tabla \textbf{A} pueden relacionarse con múltiples
		filas de la tabla \textbf{B}. Se implementa mediante una
		\textbf{tabla intermedia} (también llamada tabla de unión o pivote).
	\end{definicion}
	
	\begin{lstlisting}[style=sql, caption={Paso 1 — poblar ambas tablas padre}]
		-- Tabla padre A
		INSERT INTO estudiantes (nombre)
		VALUES ('Miguel Ríos'),
		('Sofía Castillo'),
		('Héctor Núñez');
		
		-- Tabla padre B
		INSERT INTO cursos (titulo, creditos)
		VALUES ('Bases de Datos',    6),
		('Álgebra Lineal',    4),
		('Redes de Cómputo',  5);
	\end{lstlisting}
	
	\begin{lstlisting}[style=sql, caption={Paso 2 — insertar en la tabla intermedia}]
		INSERT INTO inscripciones (estudiante_id, curso_id, fecha_alta)
		VALUES
		(1, 1, '2024-08-20'),   -- Miguel -> Bases de Datos
		(1, 3, '2024-08-20'),   -- Miguel -> Redes de Cómputo
		(2, 1, '2024-08-21'),   -- Sofía -> Bases de Datos
		(2, 2, '2024-08-21'),   -- Sofía -> Álgebra Lineal
		(3, 1, '2024-08-22'),   -- Héctor -> Bases de Datos
		(3, 2, '2024-08-22'),   -- Héctor -> Álgebra Lineal
		(3, 3, '2024-08-22');   -- Héctor -> Redes de Cómputo
	\end{lstlisting}
	
	%==============================================================
	% RESUMEN Y CONCLUSIONES
	%==============================================================
	
	\newpage
	\section*{Resumen General}
	
	Este material ha cubierto los aspectos fundamentales del lenguaje SQL,
	organizados en torno a tres ejes principales:
	
	\begin{enumerate}
		\item \textbf{Fundamentos teóricos} (secciones 1–4): Historia del problema
		de los datos, el modelo relacional de Codd, el nacimiento de SQL y
		la definición conceptual de una base de datos relacional.
		
		\item \textbf{Administración de esquemas} (secciones 5–6): Creación y
		gestión de bases de datos y tablas mediante DDL, incluyendo tipos
		de datos, restricciones y modificación de estructuras.
		
		\item \textbf{Operaciones sobre datos} (secciones 7–9): Consulta (SELECT),
		escritura (INSERT, UPDATE, DELETE) y manejo de relaciones entre
		tablas (JOIN, integridad referencial).
		
	\end{enumerate}
	
	El dominio de estos conceptos proporciona la base necesaria para:
	diseñar esquemas normalizados, escribir consultas eficientes, mantener
	la integridad de los datos y administrar bases de datos relacionales
	de forma profesional.
	
	\newpage
	\printbibliography
	
\end{document}